% ============================================================
\chapter{Martingales --- Temps Discret}
\label{ch:martingales_discret}
% ============================================================

Le terme \og martingale \fg vient du monde des jeux de hasard~: c'est la strat\'egie qui consiste \`a doubler sa mise apr\`es chaque perte, dans l'espoir de tout regagner d'un coup. Joseph Leo Doob a emprunt\'e ce mot pour d\'esigner un concept math\'ematique beaucoup plus profond~: un processus dont la meilleure pr\'ediction du futur, \'etant donn\'e le pass\'e, est la valeur actuelle. Les martingales sont l'outil central de la th\'eorie moderne des probabilit\'es~: elles interviennent dans les th\'eor\`emes de convergence, la th\'eorie de l'int\'egration stochastique, la finance math\'ematique (les prix actualis\'es sont des martingales sous la mesure risque-neutre), et la th\'eorie de l'arr\^et optimal.

\section{D\'efinitions}

\begin{definition}[Martingale en temps discret]
Soit $(\Omega, \mathcal{F}, (\mathcal{F}_n)_{n \geq 0}, \PP)$ un espace de
probabilit\'e filtr\'e. Un processus adapt\'e $(M_n)_{n \geq 0}$ est une
\emph{martingale} si~:
\begin{enumerate}
  \item $\E[\abs{M_n}] < \infty$ pour tout $n \geq 0$,
  \item $\E[M_{n+1} \mid \mathcal{F}_n] = M_n$ p.s.\ pour tout $n \geq 0$.
\end{enumerate}
On d\'efinit de mani\`ere analogue~:
\begin{itemize}
  \item \emph{sur-martingale}~: $\E[M_{n+1} \mid \mathcal{F}_n] \leq M_n$ p.s.
  \item \emph{sous-martingale}~: $\E[M_{n+1} \mid \mathcal{F}_n] \geq M_n$ p.s.
\end{itemize}
\end{definition}

\begin{remarque}
Par la propri\'et\'e de tour, la condition de martingale est \'equivalente \`a~:
pour tous $m \leq n$, $\E[M_n \mid \mathcal{F}_m] = M_m$ p.s.
En particulier, $\E[M_n] = \E[M_0]$ pour tout $n$.
\end{remarque}

\begin{exemple}
\begin{enumerate}
  \item \textbf{Marche al\'eatoire}~: si $(Y_k)_{k \geq 1}$ sont i.i.d.\ centr\'ees
    et $S_n = \sum_{k=1}^n Y_k$, alors $(S_n)$ est une martingale par rapport \`a
    $\mathcal{F}_n = \sigma(Y_1, \ldots, Y_n)$.
  \item \textbf{Transform\'ee de martingale}~: si $(M_n)$ est une martingale et $(H_n)$
    est pr\'evisible born\'e, alors $(H \cdot M)_n = \sum_{k=1}^n H_k (M_k - M_{k-1})$
    est une martingale.
  \item \textbf{Martingale exponentielle}~: si $(Y_k)$ sont i.i.d.\ et
    $\psi(\theta) = \log \E[e^{\theta Y_1}] < \infty$, alors
    $\exp\bigl(\theta S_n - n\psi(\theta)\bigr)$ est une martingale.
  \item \textbf{Martingale de L\'evy}~: si $X \in L^1$ et $\mathcal{F}_n$ est une
    filtration, alors $M_n = \E[X \mid \mathcal{F}_n]$ est une martingale.
\end{enumerate}
\end{exemple}

\section{Transform\'ee de martingale}

\begin{definition}[Transform\'ee de martingale]
Soit $(M_n)_{n \geq 0}$ une martingale et $(H_n)_{n \geq 1}$ un processus
pr\'evisible. La \emph{transform\'ee de martingale} est
\[
  (H \cdot M)_n = \sum_{k=1}^{n} H_k (M_k - M_{k-1}), \quad (H \cdot M)_0 = 0.
\]
\end{definition}

\begin{theoreme}[Stabilit\'e]
Si $(M_n)$ est une martingale (resp.\ sur-martingale) et $(H_n)$ est pr\'evisible
avec $0 \leq H_n \leq C$ pour une constante $C$, alors $(H \cdot M)_n$ est une
martingale (resp.\ sur-martingale).
\end{theoreme}

\begin{intuition}[title=\textbf{Strat\'egie de jeu}]
La transform\'ee de martingale mod\'elise les gains cumul\'es d'un joueur qui mise
$H_n$ \`a l'\'etape $n$ dans un jeu \'equitable $(M_n)$. La pr\'evisibilit\'e signifie
que la mise est d\'ecid\'ee avant de conna\^itre le r\'esultat. Le th\'eor\`eme de
stabilit\'e dit qu'aucune strat\'egie de mise born\'ee ne peut rendre un jeu
\'equitable favorable.
\end{intuition}

\section{In\'egalit\'es fondamentales}

\begin{theoreme}[In\'egalit\'e maximale de Doob]
\label{thm:doob_maximal}
Soit $(M_n)_{0 \leq n \leq N}$ une sous-martingale. Alors pour tout $\lambda > 0$~:
\[
  \lambda \, \PP\Bigl(\max_{0 \leq n \leq N} M_n \geq \lambda\Bigr)
  \leq \E[M_N^+].
\]
\end{theoreme}

\begin{theoreme}[In\'egalit\'e $L^p$ de Doob]
\label{thm:doob_Lp}
Soit $(M_n)_{0 \leq n \leq N}$ une martingale et $p > 1$. Alors
\[
  \E\Bigl[\max_{0 \leq n \leq N} \abs{M_n}^p\Bigr]
  \leq \Bigl(\frac{p}{p-1}\Bigr)^p \E[\abs{M_N}^p].
\]
\end{theoreme}

\begin{theoreme}[In\'egalit\'e de Doob pour les martingales $L^2$]
\label{thm:doob_L2}
Pour une martingale $(M_n)$ de carr\'e int\'egrable~:
\[
  \E\Bigl[\max_{0 \leq n \leq N} M_n^2\Bigr] \leq 4 \E[M_N^2].
\]
\end{theoreme}

\section{D\'ecomposition de Doob}

\begin{theoreme}[D\'ecomposition de Doob]
\label{thm:doob_decomp}
Tout processus adapt\'e int\'egrable $(X_n)_{n \geq 0}$ admet une d\'ecomposition
unique~:
\[
  X_n = M_n + A_n,
\]
o\`u $(M_n)$ est une martingale et $(A_n)$ est un processus pr\'evisible avec
$A_0 = 0$. Le processus $A$ est donn\'e par
\[
  A_n = \sum_{k=1}^{n} \bigl(\E[X_k \mid \mathcal{F}_{k-1}] - X_{k-1}\bigr).
\]
Si $(X_n)$ est une sous-martingale, alors $(A_n)$ est croissant.
\end{theoreme}

\section{Th\'eor\`emes de convergence}

\begin{theoreme}[Convergence des martingales --- Th\'eor\`eme de Doob]
\label{thm:doob_convergence}
Soit $(M_n)_{n \geq 0}$ une sous-martingale telle que $\sup_n \E[M_n^+] < \infty$.
Alors il existe une v.a.\ $M_{\infty} \in L^1$ telle que $M_n \to M_{\infty}$
p.s.
\end{theoreme}

\begin{proof}[Esquisse de la preuve]
On utilise le \emph{lemme d'upcrossing} de Doob~: si $U_N([a,b])$ d\'esigne le
nombre de remont\'ees de $(M_n)_{0 \leq n \leq N}$ de $a$ \`a $b$ (avec $a < b$), alors
\[
  (b - a) \E[U_N([a,b])] \leq \E[(M_N - a)^+].
\]
Sous l'hypoth\`ese $\sup_n \E[M_n^+] < \infty$, on obtient
$\E[U_{\infty}([a,b])] < \infty$ pour tous $a < b$ rationnels, ce qui implique
la convergence p.s.
\end{proof}

\begin{attention}[title=\textbf{Convergence p.s.\ vs.\ $L^1$}]
Le th\'eor\`eme de Doob donne la convergence p.s., mais \emph{pas} la convergence
$L^1$ en g\'en\'eral. Pour avoir $M_n \to M_{\infty}$ dans $L^1$, il faut une
condition suppl\'ementaire (uniforme int\'egrabilit\'e).
\end{attention}

\begin{definition}[Uniforme int\'egrabilit\'e]
Une famille $(X_i)_{i \in I}$ de v.a.\ est \emph{uniform\'ement int\'egrable} (UI) si
\[
  \lim_{a \to \infty} \sup_{i \in I} \E[\abs{X_i} \mathbf{1}_{\{\abs{X_i} > a\}}] = 0.
\]
\end{definition}

\begin{theoreme}[Convergence $L^1$]
\label{thm:L1_convergence}
Soit $(M_n)$ une martingale. Les assertions suivantes sont \'equivalentes~:
\begin{enumerate}
  \item $(M_n)$ est uniform\'ement int\'egrable,
  \item $M_n$ converge p.s.\ et dans $L^1$ vers $M_{\infty}$,
  \item il existe $X \in L^1$ tel que $M_n = \E[X \mid \mathcal{F}_n]$ pour tout $n$.
\end{enumerate}
Dans ce cas, $M_{\infty} = X$ p.s.\ et on dit que la martingale est \emph{ferm\'ee}.
\end{theoreme}

\begin{corollaire}[Convergence $L^p$ pour $p > 1$]
Si $(M_n)$ est une martingale born\'ee dans $L^p$ pour $p > 1$
(c'est-\`a-dire $\sup_n \E[\abs{M_n}^p] < \infty$), alors $(M_n)$ converge p.s.\ et
dans $L^p$.
\end{corollaire}

\section{Th\'eor\`eme d'arr\^et optionnel}

\begin{theoreme}[Th\'eor\`eme d'arr\^et optionnel --- version born\'ee]
\label{thm:arret_borne}
Soit $(M_n)_{n \geq 0}$ une martingale et $\tau$ un temps d'arr\^et \emph{born\'e}
($\tau \leq N$ p.s.\ pour un $N < \infty$). Alors
\[
  \E[M_{\tau}] = \E[M_0].
\]
\end{theoreme}

\begin{proof}
On \'ecrit $M_{\tau} = M_0 + \sum_{k=1}^N H_k (M_k - M_{k-1})$ o\`u
$H_k = \mathbf{1}_{\{\tau \geq k\}}$ est pr\'evisible (car $\{\tau \geq k\} = \{\tau \leq k-1\}^c
\in \mathcal{F}_{k-1}$). La transform\'ee de martingale \'etant une martingale, on a
$\E[M_{\tau}] = \E[M_0]$.
\end{proof}

\begin{theoreme}[Th\'eor\`eme d'arr\^et optionnel --- version g\'en\'erale]
\label{thm:arret_general}
Soit $(M_n)_{n \geq 0}$ une martingale uniform\'ement int\'egrable et $\sigma, \tau$
deux temps d'arr\^et avec $\sigma \leq \tau$ p.s. Alors
\[
  \E[M_{\tau} \mid \mathcal{F}_{\sigma}] = M_{\sigma} \quad \text{p.s.}
\]
En particulier, $\E[M_{\tau}] = \E[M_0]$.
\end{theoreme}

\begin{attention}[title=\textbf{Conditions du th\'eor\`eme d'arr\^et}]
Le th\'eor\`eme d'arr\^et optionnel ne s'applique \emph{pas} \`a un temps d'arr\^et
quelconque sans condition suppl\'ementaire. Contre-exemple classique~: la marche
al\'eatoire sym\'etrique $(S_n)$ avec $\tau = \inf\{n : S_n = 1\}$. On a
$\E[S_0] = 0$ mais $S_{\tau} = 1$, car $\tau$ n'est pas born\'e et $(S_n)$ n'est
pas UI.
\end{attention}

\section{Applications}

\subsection{Ruine du joueur revisit\'ee}

\begin{exemple}
Marche al\'eatoire $(S_n)$ avec $S_0 = i$ et barri\`eres absorbantes en $0$ et $N$.
On pose $\tau = \inf\{n : S_n \in \{0, N\}\}$.

Si $p \neq 1/2$, la martingale exponentielle $M_n = (q/p)^{S_n}$ est born\'ee
entre $(q/p)^0 = 1$ et $(q/p)^N$, donc $\tau$ est un temps d'arr\^et born\'e
pour $(M_{n \wedge \tau})$. Par le th\'eor\`eme d'arr\^et~:
\[
  (q/p)^i = \E[(q/p)^{S_\tau}] = (q/p)^0 \PP_i(\text{ruine})
  + (q/p)^N (1 - \PP_i(\text{ruine})),
\]
d'o\`u $\PP_i(\text{ruine}) = \frac{(q/p)^i - (q/p)^N}{1 - (q/p)^N}$.
\end{exemple}

\subsection{Loi du logarithme it\'er\'e --- premier pas}

\begin{proposition}
Si $(M_n)$ est une martingale de carr\'e int\'egrable avec
$\E[(M_n - M_{n-1})^2 \mid \mathcal{F}_{n-1}] = \sigma^2$ p.s., alors
$\langle M \rangle_n = n\sigma^2$ et $M_n^2 - n\sigma^2$ est une martingale.
\end{proposition}

\section{Crochet pr\'evisible}

\begin{definition}[Crochet pr\'evisible]
Soit $(M_n)$ une martingale de carr\'e int\'egrable. Le \emph{crochet pr\'evisible}
$(\langle M \rangle_n)$ est l'unique processus pr\'evisible croissant tel que
$(M_n^2 - \langle M \rangle_n)$ soit une martingale. Il est donn\'e par~:
\[
  \langle M \rangle_n = \sum_{k=1}^{n} \E[(M_k - M_{k-1})^2 \mid \mathcal{F}_{k-1}].
\]
\end{definition}

\begin{theoreme}[Convergence par le crochet]
\label{thm:convergence_crochet}
Soit $(M_n)$ une martingale de carr\'e int\'egrable. Si
$\langle M \rangle_{\infty} = \lim_{n} \langle M \rangle_n < \infty$ p.s.,
alors $(M_n)$ converge p.s.\ et dans $L^2$.
\end{theoreme}

\section{Exercices}

\begin{exercice}
Montrer que si $(M_n)$ est une martingale et $\varphi$ est convexe, alors
$(\varphi(M_n))$ est une sous-martingale (sous r\'eserve d'int\'egrabilit\'e).
\end{exercice}

\begin{exercice}
Soit $(X_n)_{n \geq 0}$ une marche al\'eatoire simple sym\'etrique sur $\Z$ partant
de $0$. Montrer que $X_n^2 - n$ est une martingale.
\end{exercice}

\begin{exercice}
Soit $\tau = \inf\{n \geq 0 : \abs{S_n} = a\}$ o\`u $(S_n)$ est la marche
al\'eatoire sym\'etrique. Utiliser le th\'eor\`eme d'arr\^et pour calculer $\E[\tau]$.
\end{exercice}

\begin{exercice}
D\'emontrer l'in\'egalit\'e maximale de Doob (th\'eor\`eme~\ref{thm:doob_maximal}).
\emph{Indication~: utiliser $M_N^+ \geq M_{\sigma}^+$ o\`u
$\sigma = \inf\{n : M_n \geq \lambda\} \wedge N$.}
\end{exercice}

\begin{exercice}
Soit $(M_n)$ une martingale positive avec $M_0 = 1$. Montrer que $\PP(\sup_n M_n \geq a) \leq 1/a$
pour tout $a > 0$.
\end{exercice}

\begin{exercice}
Soit $(M_n)$ une martingale born\'ee dans $L^2$. Montrer qu'elle est uniform\'ement
int\'egrable.
\end{exercice}
